\chapter{Υλοποίηση του έξυπνου συμβολαίου}
\label{ch:contract}

Το κεφάλαιο περιγράφει το πακέτο \en{Move} που υλοποιεί το πρωτόκολλο του
Κεφαλαίου~\ref{ch:apaitiseis}. Η παρουσίαση ακολουθεί τη δομή της αντίστοιχης ενότητας της
\cite{hadjichristofi2022}: μοντέλο δεδομένων, συναρτήσεις, έλεγχοι και ανάπτυξη στο δίκτυο,
με την προσθήκη μιας ανάλυσης ασφάλειας. Ο πλήρης πηγαίος κώδικας παρατίθεται στο
Παράρτημα~\ref{app:move}· τα αποσπάσματα του κεφαλαίου αντλούνται απευθείας από τα ίδια
αρχεία και διατηρούν την αρίθμηση γραμμών τους.

\section{Δομή του πακέτου}
\label{sec:package}

Το πακέτο \code{nft\_barter} αποτελείται από δύο \en{modules}, που χωρίζονται με κριτήριο
ποιες αναλλοίωτες συνθήκες προστατεύει το καθένα: το \code{marketplace\_nft} (85 γραμμές)
ορίζει τον τύπο \en{NFT} της εφαρμογής και τη δημιουργία του, ενώ το \code{barter}
(238 γραμμές) ορίζει την προσφορά, τη δέσμευση και τον κύκλο ζωής της. Το πρώτο δεν
γνωρίζει τίποτε για προσφορές και το δεύτερο τίποτε για τον συγκεκριμένο τύπο \en{NFT}.
Η σουίτα ελέγχων (515 γραμμές) ακολουθεί την ίδια διαίρεση. Η δομή του αποθετηρίου είναι η
ακόλουθη:

\begin{lstlisting}[numbers=none,caption={Δομή του αποθετηρίου},label={lst:tree}]
marketplace/
  Move.toml                       # ορισμός πακέτου, έκδοση γλώσσας 2024
  Published.toml                  # διευθύνσεις δημοσίευσης ανά δίκτυο
  sources/marketplace_nft.move    # τύπος NFT, Display, mint
  sources/barter.move             # Bid, δέσμευση, accept / cancel / reject
  tests/barter_fixture.move       # σκηνικό ελέγχων και δύο ψευδο-τύποι
  tests/barter_tests.move         # 18 έλεγχοι του κύκλου ζωής
  tests/marketplace_nft_tests.move
  scripts/                        # deploy.sh, dev.sh, sync-env.sh, ...
  app/                            # η εφαρμογή ιστού (Κεφάλαιο 5)
\end{lstlisting}

Το \code{Move.toml} δηλώνει το όνομα του πακέτου, τη σταθερή έκδοση \en{2024} της γλώσσας
και τα δίκτυα στα οποία δημοσιεύεται. Δεν δηλώνονται ρητές εξαρτήσεις· η βασική βιβλιοθήκη
(\code{std}) και το πλαίσιο του \en{Sui} (\code{sui}) παρέχονται σιωπηρά από τη
σύγχρονη αλυσίδα εργαλείων. Το πακέτο δεν εξαρτάται από κανένα πακέτο τρίτων.

\subsection{Συμβάσεις κώδικα}
\label{sec:conventions}

Ο κώδικας ακολουθεί ένα σύνολο δημοσιευμένων κανόνων ύφους και αρχιτεκτονικής για
\en{Sui Move}~\cite{sui-move-skills}, που συνοψίζονται εδώ επειδή εξηγούν αρκετές από
τις επιλογές που ακολουθούν.

\begin{itemize}[nosep]
  \item \textbf{Ένα \en{module} ανά ιδιοκτήτη αναλλοίωτων.} Ο κώδικας που μοιράζεται την
        ίδια εξουσιοδότηση, τον ίδιο κύκλο ζωής και τον ίδιο λόγο αλλαγής μένει μαζί· η
        διαίρεση γίνεται μόνο σε πραγματικά όρια, όχι για να μικρύνουν τα αρχεία.
  \item \textbf{Σταθερή διάταξη ενοτήτων.} Κάθε αρχείο έχει τις ίδιες επικεφαλίδες με την
        ίδια σειρά: \en{Errors}, \en{Public Structs}, \en{Private Initialization},
        \en{Public Functions}, \en{Private Functions}, \en{Test-Only Functions}.
  \item \textbf{Περιγραφικά σφάλματα.} Οι κωδικοί σφαλμάτων δηλώνονται με την επισημείωση
        \code{\#[error(code = N)]} και φέρουν μήνυμα κειμένου (Ενότητα~\ref{sec:events-errors}).
  \item \textbf{Μικρές συναρτήσεις με ορατή σειρά φάσεων.} Κάθε δημόσια συνάρτηση είναι μια
        σύντομη ακολουθία: έλεγχοι, μετακίνηση στοιχείων, διαγραφή, συμβάν. Οι φάσεις που
        επαναλαμβάνονται ζουν σε ιδιωτικές βοηθητικές συναρτήσεις.
  \item \textbf{Σύνταξη δέκτη και μακροεντολές της βιβλιοθήκης.} Χρησιμοποιείται η σύνταξη
        \code{bid.escrow(\ldots)} αντί για \code{escrow(\&mut bid, \ldots)}, και οι
        μακροεντολές \code{do!} και \code{any!} της \code{std} αντί για χειρόγραφους βρόχους.
  \item \textbf{Έλεγχοι ως προτάσεις συμπεριφοράς.} Οι έλεγχοι ονομάζονται από αυτό που
        αποδεικνύουν και μοιράζονται ένα κοινό σκηνικό (Ενότητα~\ref{sec:tests}).
\end{itemize}

\section{Μοντέλο δεδομένων}
\label{sec:data-model}

Τα δύο \en{modules} ορίζουν τέσσερις δομές κατάστασης και πέντε δομές συμβάντων.
Το Σχήμα~\ref{fig:class} απεικονίζει τις πρώτες σε μορφή διαγράμματος κλάσεων.

\begin{figure}[ht]
\centering
\begin{tikzpicture}[
  font=\footnotesize, >=Stealth,
  cls/.style={draw, rectangle split, rectangle split parts=2, rectangle split part fill={gray!15,white}, align=left, inner sep=5pt, text width=#1},
  cls/.default=6.4cm]
  \node[cls=6.6cm] (bid) at (0,0) {
    \textbf{\code{Bid<Offered, Requested>}}\\
    {\scriptsize\emph{key} --- \code{barter}, παράμετροι \code{phantom}}
    \nodepart{second}
    \code{id: UID}\\
    \code{bidder: address}\\
    \code{offered\_count: u64}\\
    \code{offered\_sui: Balance<SUI>}\\
    \code{requested\_count: u64}\\
    \code{target\_nft\_id: Option<ID>}};
  \node[cls=4.8cm, anchor=north west] (nft) at ([xshift=1.4cm]bid.north east) {
    \textbf{\code{MarketplaceNFT}}\\
    {\scriptsize\emph{key, store} --- \code{marketplace\_nft}}
    \nodepart{second}
    \code{id: UID}\\
    \code{name: String}\\
    \code{description: String}\\
    \code{image\_url: String}\\
    \code{collection: String}\\
    \code{rarity: String}};
  \node[cls=6.6cm, anchor=north west] (key) at ([yshift=-1.2cm]bid.south west) {
    \textbf{\code{EscrowKey(u64)}}\\
    {\scriptsize\emph{copy, drop, store}}
    \nodepart{second}
    κλειδί δυναμικού πεδίου αντικειμένου\\
    τιμή: ένα αντικείμενο τύπου \code{Offered}};
  \node[cls=4.8cm, anchor=north west] (otw) at ([yshift=-0.9cm]nft.south west) {
    \textbf{\code{MARKETPLACE\_NFT}}\\
    {\scriptsize\emph{drop}}
    \nodepart{second}
    \en{one-time witness} για το \code{Publisher}};
  \draw[->] (bid.south) -- node[right, font=\scriptsize] {\en{DOF}: \code{EscrowKey(0)} \ldots{} \code{EscrowKey(n-1)}} (key.north);
\end{tikzpicture}
\caption{Οι δομές κατάστασης των δύο \en{modules}. Το \code{Bid} δεν περιέχει τα
προσφερόμενα \en{NFTs} ως πεδίο· αυτά προσαρτώνται ως δυναμικά πεδία αντικειμένων με
κλειδιά \code{EscrowKey(0)}, \code{EscrowKey(1)}, κ.ο.κ.}
\label{fig:class}
\end{figure}

\subsection{Η δομή Bid}

\lstinputlisting[language=Move,firstline=27,lastline=40,firstnumber=27,
  caption={Η δομή της προσφοράς και το κλειδί δέσμευσης},label={lst:bid}]{\bartersrc}

Η δομή \code{Bid} είναι παραμετρική ως προς δύο τύπους. Και οι δύο περιορίζονται σε
\code{key + store}, δηλαδή σε αντικείμενα που μπορούν να αποθηκευθούν και να μεταβιβασθούν
ελεύθερα, ιδιότητες που έχει κάθε τυπικό \en{NFT}. Και οι δύο είναι \code{phantom}, αφού
κανένα πεδίο δεν έχει τύπο \code{Offered} ή \code{Requested}: οι παράμετροι υπάρχουν μόνο
για να διακρίνουν, σε επίπεδο τύπων, μια προσφορά «σπαθιών για ασπίδα» από μια προσφορά
«ασπίδων για σπαθί». Ο Πίνακας~\ref{tab:bid-fields} εξηγεί τα πεδία.

\begin{table}[H]
\centering
\caption{Πεδία της δομής \code{Bid}}
\label{tab:bid-fields}
\small
\begin{tabular}{lp{10.6cm}}
\toprule
Πεδίο & Ρόλος \\
\midrule
\code{id} & Το μοναδικό αναγνωριστικό του αντικειμένου. Σε αυτό προσαρτώνται τα δυναμικά πεδία. \\
\code{bidder} & Η διεύθυνση του προσφέροντος. Χρησιμοποιείται για την εξουσιοδότηση της ακύρωσης και ως παραλήπτης σε κάθε διαδρομή επιστροφής. \\
\code{offered\_count} & Το πλήθος των δεσμευμένων \en{NFTs}, άρα και το εύρος των κλειδιών \code{EscrowKey}. \\
\code{offered\_sui} & Το δεσμευμένο ποσό \en{SUI}, ως \code{Balance} (ενδεχομένως μηδενικό). \\
\code{requested\_count} & Το πλήθος αντικειμένων τύπου \code{Requested} που ζητούνται. \\
\code{target\_nft\_id} & Προαιρετικά, το \en{ID} του συγκεκριμένου \en{NFT} που ζητείται. \\
\bottomrule
\end{tabular}
\end{table}

Το \code{Bid} έχει μόνο την ικανότητα \code{key}: προορίζεται αποκλειστικά για κοινόχρηστο
αντικείμενο κορυφαίου επιπέδου και κοινοποιείται με την \code{transfer::share\_object}. Η
απουσία του \code{store} σημαίνει ότι κανένα άλλο πακέτο δεν μπορεί να το περιτυλίξει ή να
το μεταβιβάσει· ο μόνος τρόπος να φύγει μια προσφορά από την αλυσίδα είναι οι τρεις
συναρτήσεις τερματισμού του \en{module}. Αξίζει επίσης να σημειωθεί τι \emph{δεν}
αποθηκεύεται: τα ονόματα των δύο τύπων. Αυτά περιέχονται ήδη στον πλήρη τύπο του
αντικειμένου, όπως τον αναφέρει κάθε κόμβος, και η αντιγραφή τους σε πεδία θα ήταν
πλεονασμός. Το κλειδί \code{EscrowKey} είναι μια δομή-περιτύλιγμα ενός ακεραίου. Η χρήση
ιδιαίτερου τύπου αντί για γυμνό \code{u64} εξασφαλίζει ότι τα πεδία δέσμευσης δεν θα
συγκρουσθούν ποτέ με δυναμικά πεδία άλλης προέλευσης, διότι το κλειδί ενός δυναμικού
πεδίου περιλαμβάνει και τον τύπο του.

\subsection{Ο τύπος NFT της εφαρμογής}

\lstinputlisting[language=Move,firstline=10,lastline=27,firstnumber=10,
  caption={Το \en{one-time witness}, ο τύπος \code{MarketplaceNFT} και το συμβάν δημιουργίας},label={lst:nft}]{\nftsrc}

Ο \code{MarketplaceNFT} είναι ο τύπος που δημιουργεί η εφαρμογή επίδειξης. Η συλλογή και
η σπανιότητα είναι πεδία κειμένου, ώστε ένας μόνο τύπος να καλύπτει και τις τρεις συλλογές
της εφαρμογής. Πρέπει να τονισθεί ότι το πρωτόκολλο προσφορών \emph{δεν} εξαρτάται από
αυτόν τον τύπο: το \code{barter} δεν τον εισάγει καν, και οι έλεγχοι της
Ενότητας~\ref{sec:tests} ασκούν τις προσφορές με δύο άσχετους μεταξύ τους τύπους.

\section{Αρχικοποίηση και Display}
\label{sec:init}

\lstinputlisting[language=Move,firstline=31,lastline=46,firstnumber=31,
  caption={Η συνάρτηση \code{init} του \code{marketplace\_nft}},label={lst:init}]{\nftsrc}

Η \code{init} εκτελείται μία φορά, κατά τη δημοσίευση. Παραλαμβάνει το \en{one-time
witness} \code{MARKETPLACE\_NFT} (Ενότητα~\ref{sec:otw}) και με αυτό διεκδικεί το
αντικείμενο \code{Publisher}. Έπειτα ορίζει το \en{Display} του \code{MarketplaceNFT}:
πέντε πρότυπα που αντιστοιχίζουν τα τυπικά ονόματα προβολής στα ομώνυμα πεδία της δομής,
με τις δύο λίστες να παράγονται από τη μακροεντολή \code{map!} πάνω σε κυριολεκτικά
διανύσματα. Τέλος, μεταβιβάζει το \code{Publisher} και το \code{Display} στον εκδότη, ώστε
να μπορεί να τροποποιήσει τα πρότυπα στο μέλλον. Κανένα από τα δύο αντικείμενα δεν
παρέχει δικαιώματα επί των προσφορών· ο εκδότης του πακέτου δεν έχει κανένα προνόμιο στο
πρωτόκολλο (ΜΛΑ4). Το \code{barter} δεν έχει \code{init}: δεν χρειάζεται καμία αρχική
κατάσταση.

\section{Συναρτήσεις}
\label{sec:functions}

Ο Πίνακας~\ref{tab:functions} συνοψίζει τις δημόσιες συναρτήσεις. Είναι μόνο πέντε, όλες
\code{public}, άρα καλούνται τόσο απευθείας από \en{PTB} όσο και από άλλα πακέτα. Δεν
υπάρχουν συναρτήσεις ανάγνωσης πεδίων: οι εφαρμογές διαβάζουν τα αντικείμενα απευθείας από
τον κόμβο, και ένα ελάχιστο δημόσιο \en{ABI} είναι ευκολότερο να διατηρηθεί συμβατό σε
μελλοντικές αναβαθμίσεις.

\begin{table}[H]
\centering
\caption{Δημόσιες συναρτήσεις του πακέτου}
\label{tab:functions}
\small
\begin{tabular}{lp{4.3cm}p{6.4cm}}
\toprule
Συνάρτηση & Ορίσματα & Αποτέλεσμα \\
\midrule
\multicolumn{3}{l}{\emph{\en{module} \code{marketplace\_nft}}} \\
\code{mint} & όνομα, περιγραφή, εικόνα, συλλογή, σπανιότητα & \emph{Επιστρέφει} το νέο \code{MarketplaceNFT}· συμβάν \code{NftMinted}. \\
\addlinespace[2pt]
\multicolumn{3}{l}{\emph{\en{module} \code{barter}}} \\
\code{place\_bid<O,R>} & \code{vector<O>}, \code{Coin<SUI>}, πλήθος ζητουμένων, \code{Option<ID>} & Νέο κοινόχρηστο \code{Bid<O,R>} με τα στοιχεία δεσμευμένα· \code{BidPlaced}. \\
\code{accept\_bid<O,R>} & \code{Bid<O,R>}, \code{vector<R>} & Ανταλλαγή και διαγραφή του \code{Bid}· \code{BidAccepted}. \\
\code{cancel\_bid<O,R>} & \code{Bid<O,R>} & Επιστροφή στον προσφέροντα και διαγραφή· \code{BidCancelled}. \\
\code{reject\_bid<O,R>} & \code{Bid<O,R>}, \code{\&R} & Επιστροφή στον προσφέροντα και διαγραφή· \code{BidRejected}. \\
\bottomrule
\end{tabular}
\end{table}

\subsection{Δημιουργία NFT}

\lstinputlisting[language=Move,firstline=50,lastline=72,firstnumber=50,
  caption={Η συνάρτηση \code{mint}},label={lst:mint}]{\nftsrc}

Η \code{mint} είναι ανοικτή σε όλους και δεν επικυρώνει τις τιμές συλλογής και
σπανιότητας. Η επιλογή είναι συνειδητή για μια εφαρμογή επίδειξης, στην οποία ο τύπος
\en{NFT} υπάρχει για να τροφοδοτεί το πρωτόκολλο προσφορών με υλικό. Σε πραγματική συλλογή
η δημιουργία θα περιοριζόταν με κάποιο αντικείμενο-δικαίωμα (\en{capability}).

Η συνάρτηση \emph{επιστρέφει} το αντικείμενο αντί να το μεταβιβάσει στον καλούντα. Η
διαφορά είναι ουσιώδης για τη σύνθεση: μέσα σε ένα \en{PTB}, το αποτέλεσμα της \code{mint}
μπορεί να τροφοδοτήσει απευθείας την \code{place\_bid}, ώστε δημιουργία και προσφορά να
γίνουν σε μία συναλλαγή, ή να παραδοθεί στον αποστολέα με την εντολή
\code{TransferObjects}. Ο έλεγχος~3 του Πίνακα~\ref{tab:tests-nft} ασκεί ακριβώς αυτή τη
ροή. Το \en{PTB} είναι υποχρεωμένο να καταναλώσει την επιστρεφόμενη τιμή, αφού ο τύπος
δεν έχει \code{drop}· μια συναλλαγή που θα την άφηνε αχρησιμοποίητη απορρίπτεται πριν την
εκτέλεση.

\subsection{Υποβολή προσφοράς}
\label{sec:place-bid}

\lstinputlisting[language=Move,firstline=70,lastline=102,firstnumber=70,
  caption={Η συνάρτηση \code{place\_bid}},label={lst:place}]{\bartersrc}

Η συνάρτηση εκτελεί τέσσερα βήματα.

\begin{enumerate}[nosep]
  \item \textbf{Έλεγχοι εγκυρότητας} (γραμμές 78 και 79). Η προσφορά πρέπει να περιέχει
        τουλάχιστον ένα \en{NFT} ή θετικό ποσό (\code{EEmptyBid}) και να ζητά τουλάχιστον
        ένα αντικείμενο (\code{EInvalidRequestedCount}). Η σειρά των δύο ελέγχων είναι
        μέρος της διεπαφής και επαληθεύεται από δικό της έλεγχο.
  \item \textbf{Κατασκευή και δέσμευση} (81 έως 89). Το νόμισμα μετατρέπεται σε
        \code{Balance} και αποθηκεύεται στο πεδίο \code{offered\_sui}. Έπειτα η ιδιωτική
        \code{escrow} (Απόσπασμα~\ref{lst:helpers}) προσαρτά κάθε στοιχείο του διανύσματος
        στο \code{Bid} ως δυναμικό πεδίο αντικειμένου με κλειδί \code{EscrowKey(i)}.
  \item \textbf{Συμβάν} (91 έως 99). Εκπέμπεται \code{BidPlaced} με όλα τα στοιχεία που
        χρειάζεται μια εφαρμογή για να εντοπίσει την προσφορά, συμπεριλαμβανομένων των
        ονομάτων των δύο τύπων, που διαφορετικά θα χάνονταν μαζί με το αντικείμενο.
  \item \textbf{Κοινοποίηση} (101). Το αντικείμενο γίνεται κοινόχρηστο.
\end{enumerate}

Η μακροεντολή \code{do!} της γραμμής 174 αξίζει ιδιαίτερης μνείας. Ο τύπος του
διανύσματος, \code{vector<Offered>}, δεν έχει την ικανότητα \code{drop}, αφού δεν την έχει
ούτε ο \code{Offered}. Η \code{do!} καταναλώνει το διάνυσμα στοιχείο προς στοιχείο και το καταστρέφει μόνο
όταν αδειάσει· ένας χειρόγραφος βρόχος που θα άφηνε κάποιο στοιχείο εκτός δεν θα
μεταγλωττιζόταν καν. Με άλλα λόγια, το ίδιο το σύστημα τύπων εγγυάται ότι \emph{κάθε}
\en{NFT} που εισήλθε στη συνάρτηση έχει δεσμευθεί.

Αξίζει επίσης να παρατηρηθεί τι \emph{δεν} ελέγχει η συνάρτηση: ότι το στοχευόμενο
\en{NFT} υπάρχει, ή ότι είναι πράγματι τύπου \code{Requested}. Ο έλεγχος είναι περιττός σε
αυτό το σημείο. Μια προσφορά με ανύπαρκτο ή λανθασμένο στόχο απλώς δεν μπορεί να γίνει
αποδεκτή από κανέναν, και ο προσφέρων μπορεί πάντοτε να την ακυρώσει.

\subsection{Αποδοχή}
\label{sec:accept-bid}

\lstinputlisting[language=Move,firstline=104,lastline=120,firstnumber=104,
  caption={Η συνάρτηση \code{accept\_bid}},label={lst:accept}]{\bartersrc}

Η \code{accept\_bid} παραλαμβάνει το \code{Bid} \emph{κατά τιμή}, με το όνομα \code{self}
όπως ορίζει η σύμβαση δέκτη της γλώσσας. Αυτό σημαίνει ότι η συνάρτηση αναλαμβάνει την
κυριότητά του και οφείλει, πριν επιστρέψει, να το έχει αποδομήσει πλήρως. Το σώμα της είναι
μια σύνοψη πέντε φάσεων, καθεμία σε μία γραμμή, και οι λεπτομέρειες ζουν στις ιδιωτικές
συναρτήσεις του Αποσπάσματος~\ref{lst:helpers}.

\lstinputlisting[language=Move,firstline=155,lastline=205,firstnumber=155,
  caption={Οι ιδιωτικές φάσεις του κύκλου ζωής},label={lst:helpers}]{\bartersrc}

Πρώτα η \code{assert\_counter\_offer} επαληθεύει ότι ο καλών παραδίδει ακριβώς
\code{requested\_count} αντικείμενα (\code{EInvalidCounterCount}) και, αν η προσφορά είναι
στοχευμένη, ότι ένα από αυτά έχει το ζητούμενο \en{ID} (\code{ETargetNftNotProvided}), με
τη μακροεντολή \code{any!}. Ο τύπος των αντικειμένων δεν χρειάζεται έλεγχο: η παράμετρος
είναι \code{vector<Requested>} και το \code{Bid} είναι \code{Bid<Offered, Requested>},
άρα η ταύτιση έχει ήδη επιβληθεί από την εικονική μηχανή.

Έπειτα η \code{release\_escrow} αφαιρεί τα δεσμευμένα \en{NFTs} ένα προς ένα από τα
δυναμικά πεδία και τα μεταβιβάζει στον αποδέκτη, και η \code{do!} μεταβιβάζει τα
αντικείμενα του αποδέκτη στον προσφέροντα. Η \code{destroy} αποδομεί το \code{Bid},
διαγράφει το \code{UID} του και καταβάλλει το υπόλοιπο \en{SUI} στον αποδέκτη. Εδώ το
σύστημα τύπων επιβάλλει μια ακόμη εγγύηση: το πεδίο \code{offered\_sui} έχει τύπο
\code{Balance<SUI>}, που δεν έχει \code{drop}. Η αποδόμηση παράγει μια τιμή
\code{Balance} η οποία \emph{πρέπει} να καταλήξει κάπου, είτε ως νόμισμα στον παραλήπτη
είτε με ρητή καταστροφή αν είναι μηδενική, αλλιώς ο κώδικας δεν μεταγλωττίζεται. Το συμβάν
εκπέμπεται τελευταίο, όταν κάθε μετακίνηση έχει ολοκληρωθεί.

Αξίζει να προσεχθεί ένα σημείο στο οποίο ο μεταγλωττιστής \emph{δεν} βοηθά. Η διαγραφή
ενός \code{UID} δεν αποτυγχάνει αν το αντικείμενο έχει ακόμη δυναμικά πεδία· τα πεδία
αυτά απλώς καθίστανται απρόσιτα. Η ορθότητα της αποδέσμευσης εξαρτάται λοιπόν από το ότι η
\code{release\_escrow} αφαιρεί ακριβώς \code{offered\_count} πεδία και ότι τόσα ακριβώς
πρόσθεσε η \code{escrow}. Η ισότητα αυτή ισχύει εκ κατασκευής, αφού το
\code{offered\_count} υπολογίζεται από το μήκος του διανύσματος και δεν τροποποιείται
ποτέ, και επαληθεύεται από τους ελέγχους με πολλαπλά δεσμευμένα \en{NFTs}.

\subsection{Ακύρωση και απόρριψη}
\label{sec:cancel-reject}

\lstinputlisting[language=Move,firstline=122,lastline=153,firstnumber=122,
  caption={Οι συναρτήσεις \code{cancel\_bid} και \code{reject\_bid}},label={lst:cancel}]{\bartersrc}

Οι δύο συναρτήσεις έχουν το ίδιο σώμα με την αποδοχή, χωρίς τη φάση των αντικειμένων του
αποδέκτη, και διαφέρουν μόνο στην εξουσιοδότηση. Η ακύρωση επιτρέπεται μόνο στον
προσφέροντα, και ο έλεγχος είναι μια απλή σύγκριση της αποθηκευμένης διεύθυνσης με τον
αποστολέα της συναλλαγής (\code{ENotBidder}).

Η απόρριψη εφαρμόζει τον μηχανισμό εξουσιοδότησης που περιγράφηκε στην
Ενότητα~\ref{sec:decisions}. Η συνάρτηση δέχεται μια αμετάβλητη αναφορά
\code{\&Requested} και ελέγχει ότι το \en{ID} του αντικειμένου ταυτίζεται με τον στόχο της
προσφοράς (\code{ENotTargetNftOwner}), με τη συνάρτηση \code{contains} του
\code{Option}. Ο αποστολέας δεν συγκρίνεται με καμία διεύθυνση. Η απόδειξη της κυριότητας
είναι το ίδιο το γεγονός ότι ο καλών μπόρεσε να δώσει το αντικείμενο ως είσοδο στη
συναλλαγή. Το \en{NFT} δεν μετακινείται· μετά τη συναλλαγή παραμένει στον κάτοχό του. Μια
μη στοχευμένη προσφορά δεν μπορεί να απορριφθεί (\code{EBidNotTargeted}).

Μια συνέπεια που αξίζει να γίνει ρητή: η προσφορά δεσμεύεται στο \emph{αντικείμενο}, όχι
στον κάτοχό του. Αν το στοχευόμενο \en{NFT} αλλάξει χέρια όσο η προσφορά εκκρεμεί, ο νέος
κάτοχος μπορεί να την αποδεχθεί ή να την απορρίψει, και ο παλιός δεν μπορεί πια να κάνει
τίποτε από τα δύο. Ο έλεγχος~18 του Πίνακα~\ref{tab:tests} επιβεβαιώνει τη συμπεριφορά.

\section{Συμβάντα και κωδικοί σφαλμάτων}
\label{sec:events-errors}

Το πακέτο εκπέμπει πέντε συμβάντα, ένα για κάθε μεταβολή κατάστασης
(Πίνακας~\ref{tab:events}). Τα ωφέλιμα φορτία τους είναι σκόπιμα ελάχιστα: περιέχουν μόνο
όσα δεν μπορούν να ανακτηθούν από αλλού. Ο αποστολέας κάθε συναλλαγής, που είναι ο
προσφέρων στο \code{BidPlaced}, ο αποδέκτης στο \code{BidAccepted} και ο απορρίπτων στο
\code{BidRejected}, καταγράφεται ήδη στα μεταδεδομένα του συμβάντος από τον κόμβο και δεν
επαναλαμβάνεται. Τα ονόματα των τύπων περιλαμβάνονται στο \code{BidPlaced} επειδή το
αντικείμενο \code{Bid} διαγράφεται στο τέλος του κύκλου ζωής του και δεν θα υπήρχε άλλη
ανθεκτική πηγή τους. Η εφαρμογή χρησιμοποιεί τα δύο πρώτα συμβάντα για την ανακάλυψη
\en{NFTs} και προσφορών· τα υπόλοιπα τρία επιτρέπουν την ανασύσταση του πλήρους ιστορικού
από οποιονδήποτε.

\begin{table}[H]
\centering
\caption{Συμβάντα του πακέτου}
\label{tab:events}
\small
\begin{tabular}{llp{8.4cm}}
\toprule
Συμβάν & \en{Module} & Πεδία \\
\midrule
\code{NftMinted} & \code{marketplace\_nft} & \code{nft\_id} \\
\code{BidPlaced} & \code{barter} & \code{bid\_id}, \code{offered\_type}, \code{requested\_type}, \code{offered\_count}, \code{offered\_sui}, \code{requested\_count}, \code{target\_nft\_id} \\
\code{BidAccepted} & \code{barter} & \code{bid\_id} \\
\code{BidCancelled} & \code{barter} & \code{bid\_id} \\
\code{BidRejected} & \code{barter} & \code{bid\_id} \\
\bottomrule
\end{tabular}
\end{table}

Οι αποτυχίες σηματοδοτούνται με επτά κωδικούς (Πίνακας~\ref{tab:errors}), που
δηλώνονται με την επισημείωση \code{\#[error(code = N)]} και συνοδεύονται από μήνυμα
κειμένου. Ο κώδικας του σφάλματος είναι ο αριθμός που βλέπει η εφαρμογή και ο εξερευνητής,
ενώ το μήνυμα, όπως «\en{A bid must offer at least one NFT or some SUI}», ενσωματώνεται
στο μεταγλωττισμένο πακέτο και εμφανίζεται από τα εργαλεία που το υποστηρίζουν. Όλοι
εγείρονται από το \code{barter}· το \code{marketplace\_nft} δεν έχει καμία συνθήκη
αποτυχίας. Η στήλη «Συνάρτηση» δείχνει πού εγείρεται ο καθένας.

\begin{table}[H]
\centering
\caption{Κωδικοί σφαλμάτων του συμβολαίου}
\label{tab:errors}
\small
\begin{tabular}{clp{5.6cm}l}
\toprule
Κωδ. & Σταθερά & Σημασία & Συνάρτηση \\
\midrule
0 & \code{ENotBidder} & Ο καλών δεν είναι ο προσφέρων & \code{cancel\_bid} \\
1 & \code{EInvalidCounterCount} & Λάθος πλήθος παραδιδόμενων \en{NFTs} & \code{accept\_bid} \\
2 & \code{ETargetNftNotProvided} & Το στοχευόμενο \en{NFT} δεν παραδόθηκε & \code{accept\_bid} \\
3 & \code{EEmptyBid} & Προσφορά χωρίς \en{NFT} και χωρίς \en{SUI} & \code{place\_bid} \\
4 & \code{EInvalidRequestedCount} & Μηδενικό πλήθος ζητουμένων & \code{place\_bid} \\
5 & \code{EBidNotTargeted} & Απόρριψη μη στοχευμένης προσφοράς & \code{reject\_bid} \\
6 & \code{ENotTargetNftOwner} & Η αναφορά δεν είναι το στοχευόμενο \en{NFT} & \code{reject\_bid} \\
\bottomrule
\end{tabular}
\end{table}

\section{Ανάλυση ασφάλειας}
\label{sec:security}

Η ενότητα εξετάζει κατά πόσο η υλοποίηση ικανοποιεί τις μη λειτουργικές απαιτήσεις και τις
αναλλοίωτες συνθήκες του Κεφαλαίου~\ref{ch:apaitiseis}, και καταγράφει τα όρια της
ανάλυσης.

\subsection{Εξουσιοδότηση}

Ο Πίνακας~\ref{tab:authz} δείχνει ποιος μπορεί να καλέσει επιτυχώς κάθε συνάρτηση και με
ποιον μηχανισμό επιβάλλεται ο περιορισμός (ΜΛΑ3).

\begin{table}[H]
\centering
\caption{Εξουσιοδότηση ανά συνάρτηση}
\label{tab:authz}
\small
\begin{tabular}{lp{4.0cm}p{6.6cm}}
\toprule
Συνάρτηση & Ποιος & Μηχανισμός \\
\midrule
\code{place\_bid} & Οποιοσδήποτε, με δικά του στοιχεία & Κυριότητα αντικειμένων: μόνο ο κάτοχος μπορεί να τα δώσει ως όρισμα κατά τιμή. \\
\code{accept\_bid} & Ο κάτοχος του στόχου· σε ανοικτή προσφορά, όποιος κατέχει αντικείμενα τύπου \code{R} & Κυριότητα των παραδιδόμενων αντικειμένων και έλεγχος \en{ID} στόχου. \\
\code{cancel\_bid} & Ο προσφέρων & Σύγκριση \code{self.bidder} με τον αποστολέα. \\
\code{reject\_bid} & Ο κάτοχος του στόχου & Απαίτηση αναφοράς \code{\&R} στο στοχευόμενο αντικείμενο. \\
\bottomrule
\end{tabular}
\end{table}

\subsection{Διατήρηση των δεσμευμένων στοιχείων}

Η ΜΛΑ2 και οι αναλλοίωτες Α1 έως Α3 ελέγχονται εξετάζοντας κάθε διαδρομή.

Κατά την \emph{υποβολή}, τα \en{NFTs} εξέρχονται από την κυριότητα του προσφέροντος και
γίνονται παιδιά του \code{Bid}, και το \en{SUI} γίνεται πεδίο του. Από τη στιγμή που το
\code{Bid} κοινοποιείται, κανείς δεν έχει άμεση πρόσβαση στα περιεχόμενά του: τα δυναμικά
πεδία και το \code{Balance} είναι προσβάσιμα μόνο από κώδικα του \en{module} που όρισε τον
τύπο, και το \code{barter} δεν διαθέτει άλλη συνάρτηση που τα αγγίζει πέραν των τριών
συναρτήσεων τερματισμού (Α1).

Κατά τον \emph{τερματισμό}, και οι τρεις συναρτήσεις περνούν από τις ίδιες δύο ιδιωτικές
φάσεις: η \code{release\_escrow} αφαιρεί όλα τα δυναμικά πεδία και η \code{destroy}
αποδομεί το \code{Bid}, διαγράφει το \code{UID} και διαθέτει το \code{Balance}. Ο παραλήπτης είναι ο αποδέκτης στην αποδοχή και ο προσφέρων στις άλλες
δύο. Εφόσον το \code{Bid} παραλαμβάνεται κατά τιμή και διαγράφεται, δεν υπάρχει μετά τη
συναλλαγή αντικείμενο με το οποίο θα μπορούσε να επαναληφθεί η ενέργεια (Α2). Τα στοιχεία
που εξέρχονται είναι εκείνα που εισήλθαν, αφού ανακτώνται με τα ίδια κλειδιά και κανένας
κώδικας δεν μπορεί να τα αντικαταστήσει στο μεταξύ (Α3).

Σε όλες τις διαδρομές, το σύστημα τύπων της \en{Move} αποκλείει δύο ολόκληρες κατηγορίες
σφαλμάτων: την παράλειψη διάθεσης ενός \en{NFT} που βρίσκεται σε τοπική μεταβλητή ή
διάνυσμα, και την παράλειψη διάθεσης του \code{Balance}.

\subsection{Ατομικότητα και ταυτόχρονη πρόσβαση}

Η ατομικότητα (ΜΛΑ1) παρέχεται από το μοντέλο εκτέλεσης: αν οποιαδήποτε \code{assert!}
αποτύχει, ολόκληρη η συναλλαγή αναιρείται. Δεν υπάρχει κατάσταση στην οποία ο ένας έχει
παραλάβει και ο άλλος όχι.

Για την ταυτόχρονη πρόσβαση (ΜΛΑ6), ας υποτεθεί ότι ο προσφέρων υποβάλλει ακύρωση την ίδια
στιγμή που ο κάτοχος υποβάλλει αποδοχή. Και οι δύο συναλλαγές δηλώνουν το ίδιο κοινόχρηστο
αντικείμενο ως είσοδο, άρα διατάσσονται από τη συναίνεση
(Ενότητα~\ref{sec:sui-consensus}). Όποια εκτελεστεί πρώτη διαγράφει το \code{Bid}. Η
δεύτερη αποτυγχάνει επειδή η είσοδός της δεν υπάρχει πλέον, χωρίς καμία παρενέργεια πέραν
του \en{gas}. Το αποτέλεσμα είναι πάντοτε μία από τις δύο τερματικές καταστάσεις, ποτέ
κάποιο μείγμα τους. Η συμπεριφορά αυτή επιβεβαιώθηκε πειραματικά στο τοπικό δίκτυο, με δύο
λογαριασμούς να υποβάλλουν ταυτόχρονα αποδοχή της ίδιας προσφοράς
(Ενότητα~\ref{sec:scenario-results}).

\subsection{Απουσία επανεισόδου}

Οι επιθέσεις επανεισόδου (\en{reentrancy}), που έχουν προκαλέσει σημαντικές απώλειες σε
συμβόλαια του \en{Ethereum}, προϋποθέτουν ότι η μεταβίβαση ενός στοιχείου παραδίδει τον
έλεγχο σε κώδικα του παραλήπτη. Στη \en{Move} του \en{Sui} αυτό δεν συμβαίνει. Η
\code{transfer::public\_transfer} απλώς αλλάζει τον ιδιοκτήτη του αντικειμένου· δεν καλεί
κανέναν κώδικα. Η γλώσσα δεν διαθέτει δυναμική αποστολή κλήσεων, και το πακέτο δεν
καλεί συναρτήσεις τρίτων πακέτων. Η κατηγορία αυτή επιθέσεων δεν έχει εφαρμογή.

\subsection{Παρατηρήσεις και γνωστά όρια}
\label{sec:security-limits}

Η ανάλυση ανέδειξε ορισμένα σημεία που δεν αποτελούν ευπάθειες αλλά πρέπει να είναι γνωστά
σε όποιον χρησιμοποιήσει ή επεκτείνει το πρωτόκολλο.

\begin{itemize}[nosep]
  \item \textbf{Ανοικτές προσφορές.} Μια προσφορά χωρίς στόχο μπορεί να γίνει αποδεκτή από
        οποιονδήποτε παραδώσει αντικείμενα του ζητούμενου τύπου, όποια και αν είναι αυτά.
        Αυτή είναι η προβλεπόμενη σημασιολογία, και η εφαρμογή δεν τη χρησιμοποιεί.
  \item \textbf{Πολλαπλά ζητούμενα με στόχο.} Όταν \code{requested\_count} $> 1$ και
        υπάρχει στόχος, ελέγχεται μόνο ότι ο στόχος περιλαμβάνεται· τα υπόλοιπα
        αντικείμενα μπορεί να είναι οποιαδήποτε του τύπου \code{R}.
  \item \textbf{Η απόρριψη προϋποθέτει αντικείμενο διεύθυνσης.} Ο μηχανισμός της αναφοράς
        αποδεικνύει κυριότητα μόνο για αντικείμενα που ανήκουν σε διεύθυνση. Αν ο τύπος
        \code{R} επέτρεπε κοινόχρηστα ή αμετάβλητα αντικείμενα, οποιοσδήποτε θα μπορούσε
        να παρουσιάσει αναφορά σε αυτά. Η συνέπεια θα ήταν μόνο ότι τρίτος θα μπορούσε να
        προκαλέσει επιστροφή των στοιχείων στον προσφέροντα, όχι απώλεια.
  \item \textbf{Απουσία λήξης.} Μια προσφορά παραμένει ενεργή επ' αόριστον, έως ότου
        κάποιος τη διεκπεραιώσει. Τα στοιχεία του προσφέροντος δεν κινδυνεύουν, αφού μπορεί
        πάντοτε να ακυρώσει, αλλά ένας κάτοχος μπορεί να αποδεχθεί μια παλιά προσφορά που
        ο προσφέρων έχει ξεχάσει.
  \item \textbf{Ανεπιθύμητες προσφορές.} Οποιοσδήποτε μπορεί να υποβάλει προσφορές για
        οποιοδήποτε \en{NFT}. Το κόστος το φέρει ο ίδιος, αφού δεσμεύει στοιχεία και
        καταβάλλει την εγγύηση αποθήκευσης, και ο κάτοχος μπορεί να τις απορρίψει ή να τις
        αγνοήσει.
  \item \textbf{Εμβέλεια της ανάλυσης.} Η ανάλυση είναι επιχειρηματολογική και στηρίζεται
        σε ελέγχους μονάδας. Δεν έχει γίνει τυπική επαλήθευση ούτε ανεξάρτητος έλεγχος
        ασφάλειας, που θα ήταν προαπαιτούμενα για χρήση με στοιχεία πραγματικής αξίας.
\end{itemize}

\section{Έλεγχοι}
\label{sec:tests}

Η σουίτα περιλαμβάνει 21 ελέγχους μονάδας γραμμένους με το \code{sui::test\_scenario}, που
προσομοιώνει διαδοχικές συναλλαγές διαφορετικών αποστολέων. Οι 18 ασκούν το \code{barter}
με δύο τύπους που ορίζονται μόνο για τους ελέγχους, \code{Sword} και \code{Shield}, οι
οποίοι δεν έχουν καμία σχέση μεταξύ τους ούτε με τον \code{MarketplaceNFT}. Με αυτόν τον
τρόπο οι έλεγχοι επαληθεύουν ακριβώς τον ισχυρισμό του τίτλου της εργασίας: ότι η
ανταλλαγή λειτουργεί μεταξύ αντικειμένων διαφορετικών τύπων. Οι υπόλοιποι τρεις ασκούν το
\code{marketplace\_nft}.

\subsection{Το σκηνικό των ελέγχων}

Οι δύο ψευδο-τύποι, το σκηνικό και όλες οι βοηθητικές συναρτήσεις ζουν σε ξεχωριστό
\en{module} ελέγχων, το \code{barter\_fixture}, ώστε κάθε έλεγχος να δηλώνει μόνο τη δική
του διαφορά. Το σκηνικό είναι μια δομή χωρίς ικανότητες που περιτυλίγει το
\code{Scenario}· η έλλειψη \code{drop} υποχρεώνει κάθε έλεγχο να τερματίσει ρητά με
\code{end()}, άρα κανένα σενάριο δεν μένει μισοτελειωμένο. Δύο μακροεντολές,
\code{next\_tx!} και \code{ctx!}, είναι ο μόνος τρόπος προώθησης των συναλλαγών και
πρόσβασης στο πλαίσιο, και μια σειρά μεθόδων όπως \code{swords(3)}, \code{sui(50)},
\code{take\_bid()}, \code{sword\_count()} και \code{sui\_received()} εκφράζουν το κάθε
βήμα στη γλώσσα του προβλήματος. Το Απόσπασμα~\ref{lst:test-example} δείχνει έναν πλήρη
έλεγχο.

\lstinputlisting[language=Move,firstline=84,lastline=104,firstnumber=84,
  caption={Ένας έλεγχος αποδοχής με το σκηνικό},label={lst:test-example}]{\testsrc}

\subsection{Τι καλύπτεται}

Ο Πίνακας~\ref{tab:tests} παραθέτει τους ελέγχους του \code{barter}, ομαδοποιημένους όπως
στο αρχείο. Τα ονόματά τους είναι προτάσεις που περιγράφουν τη συμπεριφορά που
αποδεικνύουν, όχι τη συνάρτηση που καλούν.

\begin{table}[H]
\centering
\caption{Οι έλεγχοι του \en{module} \code{barter}}
\label{tab:tests}
\footnotesize
\begin{tabular}{rp{7.2cm}p{4.6cm}l}
\toprule
\# & Έλεγχος & Σενάριο & Αναμένεται \\
\midrule
\multicolumn{4}{l}{\emph{Υποβολή}} \\
1 & \tn{place\_\allowbreak{}bid\_\allowbreak{}escrows\_\allowbreak{}the\_\allowbreak{}offer\_\allowbreak{}in\_\allowbreak{}a\_\allowbreak{}shared\_\allowbreak{}bid} & 2 σπαθιά· τα \en{IDs} τους βρίσκονται στα δυναμικά πεδία & επιτυχία \\
2 & \tn{place\_\allowbreak{}bid\_\allowbreak{}emits\_\allowbreak{}the\_\allowbreak{}full\_\allowbreak{}offer} & όλα τα πεδία του \code{BidPlaced} & επιτυχία \\
3 & \tn{place\_\allowbreak{}bid\_\allowbreak{}rejects\_\allowbreak{}an\_\allowbreak{}empty\_\allowbreak{}offer} & χωρίς \en{NFT} και χωρίς \en{SUI} & κωδικός 3 \\
4 & \tn{place\_\allowbreak{}bid\_\allowbreak{}rejects\_\allowbreak{}a\_\allowbreak{}zero\_\allowbreak{}requested\_\allowbreak{}count} & ζητούνται 0 αντικείμενα & κωδικός 4 \\
5 & \tn{place\_\allowbreak{}bid\_\allowbreak{}checks\_\allowbreak{}the\_\allowbreak{}offer\_\allowbreak{}before\_\allowbreak{}the\_\allowbreak{}requested\_\allowbreak{}count} & και τα δύο λάθος· προτεραιότητα & κωδικός 3 \\
\multicolumn{4}{l}{\emph{Αποδοχή}} \\
6 & \tn{accept\_\allowbreak{}bid\_\allowbreak{}swaps\_\allowbreak{}the\_\allowbreak{}escrow\_\allowbreak{}for\_\allowbreak{}the\_\allowbreak{}counter\_\allowbreak{}nfts} & 3 σπαθιά για 1 ασπίδα & επιτυχία \\
7 & \tn{accept\_\allowbreak{}bid\_\allowbreak{}pays\_\allowbreak{}the\_\allowbreak{}escrowed\_\allowbreak{}sui\_\allowbreak{}to\_\allowbreak{}the\_\allowbreak{}acceptor} & 2 σπαθιά και 50 \en{MIST} & επιτυχία \\
8 & \tn{accept\_\allowbreak{}bid\_\allowbreak{}works\_\allowbreak{}for\_\allowbreak{}a\_\allowbreak{}sui\_\allowbreak{}only\_\allowbreak{}offer} & μόνο \en{SUI}, στοχευμένη & επιτυχία \\
9 & \tn{accept\_\allowbreak{}bid\_\allowbreak{}delivers\_\allowbreak{}every\_\allowbreak{}requested\_\allowbreak{}nft\_\allowbreak{}to\_\allowbreak{}the\_\allowbreak{}bidder} & ζητούνται 2, στόχος και ένα ακόμη & επιτυχία \\
10 & \tn{accept\_\allowbreak{}bid\_\allowbreak{}rejects\_\allowbreak{}the\_\allowbreak{}wrong\_\allowbreak{}number\_\allowbreak{}of\_\allowbreak{}counter\_\allowbreak{}nfts} & 2 αντί για 1 & κωδικός 1 \\
11 & \tn{accept\_\allowbreak{}bid\_\allowbreak{}rejects\_\allowbreak{}a\_\allowbreak{}counter\_\allowbreak{}offer\_\allowbreak{}without\_\allowbreak{}the\_\allowbreak{}target} & άλλη ασπίδα από τη ζητούμενη & κωδικός 2 \\
12 & \tn{accept\_\allowbreak{}bid\_\allowbreak{}checks\_\allowbreak{}the\_\allowbreak{}count\_\allowbreak{}before\_\allowbreak{}the\_\allowbreak{}target} & και τα δύο λάθος· προτεραιότητα & κωδικός 1 \\
\multicolumn{4}{l}{\emph{Ακύρωση}} \\
13 & \tn{cancel\_\allowbreak{}bid\_\allowbreak{}returns\_\allowbreak{}the\_\allowbreak{}escrow\_\allowbreak{}and\_\allowbreak{}sui\_\allowbreak{}to\_\allowbreak{}the\_\allowbreak{}bidder} & 2 σπαθιά και 30 \en{MIST} πίσω & επιτυχία \\
14 & \tn{cancel\_\allowbreak{}bid\_\allowbreak{}rejects\_\allowbreak{}anyone\_\allowbreak{}but\_\allowbreak{}the\_\allowbreak{}bidder} & ακύρωση από τρίτο & κωδικός 0 \\
\multicolumn{4}{l}{\emph{Απόρριψη}} \\
15 & \tn{reject\_\allowbreak{}bid\_\allowbreak{}returns\_\allowbreak{}the\_\allowbreak{}escrow\_\allowbreak{}to\_\allowbreak{}the\_\allowbreak{}bidder} & 3 σπαθιά και 5 \en{MIST} πίσω· η ασπίδα μένει & επιτυχία \\
16 & \tn{reject\_\allowbreak{}bid\_\allowbreak{}rejects\_\allowbreak{}an\_\allowbreak{}untargeted\_\allowbreak{}bid} & απόρριψη ανοικτής προσφοράς & κωδικός 5 \\
17 & \tn{reject\_\allowbreak{}bid\_\allowbreak{}rejects\_\allowbreak{}a\_\allowbreak{}different\_\allowbreak{}nft\_\allowbreak{}than\_\allowbreak{}the\_\allowbreak{}target} & αναφορά σε άλλη ασπίδα & κωδικός 6 \\
\multicolumn{4}{l}{\emph{Κύκλος ζωής με αλλαγή κατόχου}} \\
18 & \tn{a\_\allowbreak{}bid\_\allowbreak{}follows\_\allowbreak{}the\_\allowbreak{}target\_\allowbreak{}nft\_\allowbreak{}to\_\allowbreak{}its\_\allowbreak{}new\_\allowbreak{}owner} & ο στόχος μεταβιβάζεται· ο νέος κάτοχος αποδέχεται & επιτυχία \\
\bottomrule
\end{tabular}
\end{table}

\begin{table}[H]
\centering
\caption{Οι έλεγχοι του \en{module} \tn{marketplace\_\allowbreak{}nft}}
\label{tab:tests-nft}
\footnotesize
\begin{tabular}{rp{7.2cm}p{5.4cm}}
\toprule
\# & Έλεγχος & Σενάριο \\
\midrule
1 & \tn{mint\_\allowbreak{}returns\_\allowbreak{}an\_\allowbreak{}nft\_\allowbreak{}with\_\allowbreak{}the\_\allowbreak{}given\_\allowbreak{}metadata} & και τα πέντε πεδία, συν το \code{NftMinted} \\
2 & \tn{init\_\allowbreak{}hands\_\allowbreak{}the\_\allowbreak{}publisher\_\allowbreak{}and\_\allowbreak{}display\_\allowbreak{}to\_\allowbreak{}the\_\allowbreak{}deployer} & \code{Publisher} και \code{Display} στον εκδότη \\
3 & \tn{two\_\allowbreak{}minted\_\allowbreak{}nfts\_\allowbreak{}can\_\allowbreak{}be\_\allowbreak{}bartered\_\allowbreak{}for\_\allowbreak{}one\_\allowbreak{}in\_\allowbreak{}a\_\allowbreak{}single\_\allowbreak{}ptb\_\allowbreak{}style\_\allowbreak{}flow} & \tn{mint} και \tn{place\_\allowbreak{}bid} στην ίδια συναλλαγή, έπειτα αποδοχή \\
\bottomrule
\end{tabular}
\end{table}

Δώδεκα έλεγχοι καλύπτουν επιτυχείς διαδρομές και εννέα αποτυχίες. Κάθε αποτυχία
δηλώνεται με \code{\#[expected\_failure(abort\_code = \ldots, location = \ldots)]}, που
απαιτεί να συμβεί με τον συγκεκριμένο κωδικό και στο συγκεκριμένο \en{module}. Και οι
επτά κωδικοί σφαλμάτων έχουν τον δικό τους έλεγχο. Δύο έλεγχοι (5 και 12) παρέχουν
ταυτόχρονα δύο άκυρες εισόδους και επαληθεύουν ποιος έλεγχος εγείρεται πρώτος, αφού η
σειρά των ελέγχων είναι μέρος της συμπεριφοράς που βλέπουν οι εφαρμογές. Οι έλεγχοι
συμβάντων δεν μετρούν απλώς τα συμβάντα· ο έλεγχος~2 διαβάζει και τα επτά πεδία του
\code{BidPlaced}, και οι έλεγχοι 6, 13 και 15 επαληθεύουν ότι το \code{bid\_id} του
τερματικού συμβάντος είναι το αναγνωριστικό της προσφοράς που μόλις διαγράφηκε. Η σουίτα
εκτελείται με την εντολή \code{sui move test}, ολοκληρώνεται επιτυχώς, και η μεταγλώττιση
με \code{sui move build -{}-lint} δεν παράγει καμία προειδοποίηση.

\section{Ανάπτυξη στο δίκτυο}
\label{sec:deployment}

Το πακέτο δημοσιεύεται με την εντολή \code{sui client publish}. Η δημοσίευση δημιουργεί
τέσσερα αντικείμενα: το ίδιο το πακέτο, που είναι αμετάβλητο· το \code{Publisher} και το
\code{Display} που δημιούργησε η \code{init}· και το \code{UpgradeCap}, που επιτρέπει στον
εκδότη να δημοσιεύσει νεότερη έκδοση. Η αλυσίδα εργαλείων καταγράφει τη διεύθυνση του
πακέτου ανά δίκτυο στο αρχείο \code{Published.toml}.

\subsection{Ροή ανάπτυξης ανά δίκτυο}
\label{sec:dev-flow}

Η ανάπτυξη ακολούθησε τα δίκτυα της Ενότητας~\ref{sec:sui-networks} με σειρά αυξανόμενης
μονιμότητας, ώστε κάθε στάδιο να κοστίζει όσο το δυνατόν λιγότερο σε χρόνο και σε
εξάρτηση από εξωτερικούς παράγοντες.

\begin{enumerate}[nosep]
  \item \textbf{Έλεγχοι μονάδας, χωρίς δίκτυο.} Κάθε αλλαγή στο συμβόλαιο περνά πρώτα από
        τη σουίτα της Ενότητας~\ref{sec:tests}, που εκτελείται σε δευτερόλεπτα και δεν
        χρειάζεται κόμβο.
  \item \textbf{\en{Localnet}: το κύριο περιβάλλον εργασίας.} Ολόκληρη η ανάπτυξη του
        συμβολαίου και της εφαρμογής, τα σενάρια επαλήθευσης μέσω \en{SDK}, η
        αυτοματοποιημένη πλοήγηση της διεπαφής, οι μετρήσεις \en{gas} και τα στιγμιότυπα
        του Κεφαλαίου~\ref{ch:epideixi} έγιναν σε τοπικό δίκτυο. Οι λόγοι είναι
        πρακτικοί: μηδενική καθυστέρηση, απεριόριστο \en{SUI} από το τοπικό \en{faucet},
        δυνατότητα να ξεκινά κανείς από καθαρή αλυσίδα όποτε θέλει, και ανεξαρτησία από
        τη διαθεσιμότητα δημόσιων υπηρεσιών. Το τίμημα είναι ότι κάθε επανεκκίνηση αλλάζει
        τη διεύθυνση του πακέτου, και ότι τα πορτοφόλια επέκτασης δεν βλέπουν το τοπικό
        δίκτυο, οπότε χρησιμοποιείται το προσωρινό πορτοφόλι της εφαρμογής.
  \item \textbf{\en{Devnet}: πρώτη δημόσια δοκιμή.} Μια πρώιμη έκδοση του πακέτου
        δημοσιεύθηκε στο \en{devnet}, για να επαληθευθεί η διαδικασία δημοσίευσης σε
        δημόσιο δίκτυο και η εμφάνιση των αντικειμένων στους δημόσιους εξερευνητές. Οι
        εβδομαδιαίες διαγραφές του δικτύου το καθιστούν ακατάλληλο για οτιδήποτε
        μονιμότερο, και δεν χρησιμοποιήθηκε περαιτέρω.
  \item \textbf{\en{Testnet}: η τελική έκδοση.} Το τελικό πακέτο δημοσιεύεται στο
        \en{testnet}, και η εφαρμογή διαμορφώνεται με αυτό ως προεπιλεγμένο δίκτυο. Εκεί
        το σύστημα χρησιμοποιείται όπως θα χρησιμοποιούνταν στην παραγωγή: με πραγματικό
        πορτοφόλι επέκτασης, με δημόσιο εξερευνητή για κάθε συναλλαγή, και με \en{SUI}
        από το δημόσιο \en{faucet}, που έχει όριο ρυθμού και υπενθυμίζει ότι το \en{gas}
        δεν είναι δωρεάν. Το \en{mainnet} μένει εκτός εμβέλειας: μια εφαρμογή επίδειξης
        χωρίς ανεξάρτητο έλεγχο ασφάλειας δεν πρέπει να διαχειρίζεται στοιχεία
        πραγματικής αξίας (Ενότητα~\ref{sec:security-limits}).
\end{enumerate}

Η εφαρμογή δεν χρειάζεται αλλαγή κώδικα για τη μετάβαση από το ένα δίκτυο στο άλλο: η
λίστα των δικτύων, ο πελάτης ανά δίκτυο, ο εξερευνητής και το \en{faucet} επιλέγονται από
το τρέχον δίκτυο του \en{dApp Kit} (Κεφάλαιο~\ref{ch:dapp}), και η διεύθυνση του πακέτου
ανά δίκτυο διαβάζεται από το \code{Published.toml}.

\subsection{Το τοπικό δίκτυο στην πράξη}

Το τοπικό δίκτυο εκκινείται με την εντολή
\code{sui start -{}-force-regenesis -{}-with-faucet}, που δημιουργεί κάθε φορά μια νέα,
άδεια αλυσίδα με δικό της αναγνωριστικό. Αυτό προκάλεσε μια πρακτική δυσκολία: η αλυσίδα εργαλείων «καρφώνει» το αναγνωριστικό αλυσίδας στο
\code{Move.toml} και τη διεύθυνση δημοσίευσης στο \code{Published.toml}, και μετά από κάθε
επανεκκίνηση αρνείται είτε να μεταγλωττίσει, λόγω ασυμφωνίας αναγνωριστικού, είτε να
δημοσιεύσει, θεωρώντας το πακέτο ήδη δημοσιευμένο. Η λύση ήταν τέσσερα μικρά σενάρια
φλοιού:

\begin{description}[leftmargin=1.2em,itemsep=1pt]
  \item[\code{reset-localnet-pins.sh}] Διαβάζει το αναγνωριστικό της τρέχουσας αλυσίδας,
    το γράφει στο \code{Move.toml} και αφαιρεί την παλιά εγγραφή του τοπικού δικτύου από
    το \code{Published.toml}.
  \item[\code{deploy.sh}] Εκτελεί το προηγούμενο, δημοσιεύει το πακέτο και καλεί το
    επόμενο.
  \item[\code{sync-env.sh}] Εξάγει τη νέα διεύθυνση του πακέτου από το
    \code{Published.toml} και τη γράφει στο αρχείο περιβάλλοντος της εφαρμογής ιστού.
  \item[\code{dev.sh}] Εκκινεί το τοπικό δίκτυο, περιμένει να γίνει διαθέσιμο, εκτελεί το
    \code{deploy.sh} και εκκινεί την εφαρμογή.
\end{description}

Με αυτά, ολόκληρο το περιβάλλον ανάπτυξης στήνεται με μία εντολή και η εφαρμογή γνωρίζει
πάντοτε τη σωστή διεύθυνση του πακέτου χωρίς χειροκίνητη παρέμβαση.
